% Trascrizione quiz Paola Suria: intervallo pagine del PDF.
% Ogni risposta e spiegazione e` conservata nei commenti % SOLUZIONE,
% per non alterare il parser del sito esistente.
% ATTENZIONE: verificare le voci segnalate come da verificare.

\par\noindent\hfill{\color{cyan!65!black}\scriptsize Pag. originale 21/521}\par\smallskip

% SIMULAZIONE / QUESITO 17 (pagina 21)
\begin{quiz}
Il polinomio di Taylor di II grado, con centro $x=1$, della funzione $f(x)=e^{-x^2/3}$ \`e:
\begin{enumerate}[label=\alph*),leftmargin=1.3cm,itemsep=3pt,topsep=5pt]
\item $\frac13e^{-1/3}(3-4(x-1)+\frac53(x-1)^2)$
\item $\frac13e^{-1/3}(3-2(x-1)-\frac13(x-1)^2)$
\item $\frac13e^{-1/3}(9-2(x-1)-\frac79(x-1)^2)$
\item $\frac13e^{-1/3}(9-4(x-1)-\frac19(x-1)^2)$
\item $\frac13e^{1/3}(9-4(x-1)-\frac19(x-1)^2)$
\end{enumerate}
\end{quiz}
% SOLUZIONE: b
% SPIEGAZIONE: $f(1)=e^{-1/3}$, $f\prime(1)=-2e^{-1/3}/3$, $f\prime\prime(1)=-2e^{-1/3}/9$, dunque il coefficiente quadratico \`e $-e^{-1/3}/9$.

% SIMULAZIONE / QUESITO 18 (pagina 21)
\begin{quiz}
La funzione $f:\mathbb R\setminus\{-1\}\to\mathbb R$, $f(x)=\dfrac{e^{-3x^2}}{4x+4}$ \`e crescente in
\begin{enumerate}[label=\alph*),leftmargin=1.3cm,itemsep=3pt,topsep=5pt]
\item $(-\infty,-1)$
\item $(-\frac1{2\sqrt3},1+\frac1{2\sqrt3})$
\item $(-\frac12-\frac1{2\sqrt3},-\frac12+\frac1{2\sqrt3})$
\item $(-1,+\infty)$
\item $(-\frac12-\frac1{2\sqrt3},-1)\cup(-1,-\frac12+\frac1{2\sqrt3})$
\end{enumerate}
\end{quiz}
% SOLUZIONE: c
% SPIEGAZIONE: $f\prime$ ha segno di $-6x(x+1)-1$; \`e positivo tra le radici $(-3\pm\sqrt3)/6$. Attenzione: l\'intervallo attraversa $x=-1$? No, entrambe le radici sono maggiori di $-1$.

% SIMULAZIONE / QUESITO 19 (pagina 21)
\begin{quiz}
Quali sono i valori massimo e minimo della funzione $f(x)=x^3-6x^2+9x+1$ nell\'intervallo $[-2,2]$?
\begin{enumerate}[label=\alph*),leftmargin=1.3cm,itemsep=3pt,topsep=5pt]
\item max $=1$, min $=-3$
\item max $=6$, min $=-21$
\item max $=23$, min $=-4$
\item max $=5$, min $=-49$
\item max $=51$, min $=-3$
\end{enumerate}
\end{quiz}
% SOLUZIONE: d
% SPIEGAZIONE: $f\prime=3(x-1)(x-3)$: nell\'intervallo il punto critico \`e $1$. $f(-2)=-49$, $f(1)=5$, $f(2)=3$.

% SIMULAZIONE / QUESITO 20 (pagina 21)
\begin{quiz}
Lo sviluppo di Mac Laurin di ordine 4 della funzione $f(x)=e^{\cosh x}$ \`e:
\begin{enumerate}[label=\alph*),leftmargin=1.3cm,itemsep=3pt,topsep=5pt]
\item $(1+\frac{x^2}2+\frac{x^4}{4!})+\frac12(1+\frac{x^2}2+\frac{x^4}{4!})^2+\frac1{3!}(1+\frac{x^2}2+\frac{x^4}{4!})^3+\frac1{4!}(1+\frac{x^2}2+\frac{x^4}{4!})^4+o(x^4)$
\item $e(1+\frac12x^2+\frac1{4!}x^4)+o(x^4)$
\item $e(1+\frac12x^2+\frac16x^4)+o(x^4)$
\item $1+\frac12x^2+\frac1{4!}x^4+o(x^4)$
\item $1+\frac12x^2+\frac16x^4+o(x^4)$
\end{enumerate}
\end{quiz}
% SOLUZIONE: c
% SPIEGAZIONE: $\cosh x=1+x^2/2+x^4/24+o(x^4)$; $e^{\cosh x}=e[1+x^2/2+(1/24+1/8)x^4]+o(x^4)$.

\par\noindent\hfill{\color{cyan!65!black}\scriptsize Pag. originale 23/521}\par\smallskip

% SIMULAZIONE / QUESITO 1 (pagina 23)
\begin{quiz}
Sia $z$ un numero complesso diverso da zero. Quale delle seguenti affermazioni \`e FALSA?
\begin{enumerate}[label=\alph*),leftmargin=1.3cm,itemsep=3pt,topsep=5pt]
\item $\arg(z^4)=4\arg(z)$
\item $|z^n|=|z|^n$
\item $z$ pu\`o avere meno di quattro radici quarte
\item $z+\overline z\in\mathbb R$
\item $z\overline z\in\mathbb R$
\end{enumerate}
\end{quiz}
% SOLUZIONE: c
% SPIEGAZIONE: Ogni complesso non nullo ha esattamente quattro radici quarte distinte. La formula degli argomenti si intende modulo $2\pi$.

% SIMULAZIONE / QUESITO 2 (pagina 23)
\begin{quiz}
Il dominio della funzione $f(x)=\log_2\!\left(\dfrac{\sqrt{1-x^2}}{x+4}\right)$ \`e:
\begin{enumerate}[label=\alph*),leftmargin=1.3cm,itemsep=3pt,topsep=5pt]
\item $(-\infty,-1)\cup(1,+\infty)$
\item $(-2,-1)\cup(1,2)$
\item $\mathbb R$
\item $[-1,1]$
\item $(-1,1)$
\end{enumerate}
\end{quiz}
% SOLUZIONE: e
% SPIEGAZIONE: Il radicando deve essere strettamente positivo, $1-x^2>0$, e $x+4>0$.

% SIMULAZIONE / QUESITO 3 (pagina 23)
\begin{quiz}
Sia data l\'equazione $x^{49}+4x-10=0$. Quale delle seguenti affermazioni \`e FALSA?
\begin{enumerate}[label=\alph*),leftmargin=1.3cm,itemsep=3pt,topsep=5pt]
\item l'equazione ammette una sola radice reale
\item l'equazione ammette 49 radici in $\mathbb C$
\item tra le 49 radici in $\mathbb C$, 48 sono a due a due complesse coniugate
\item la radice reale \`e positiva
\item le radici dell'equazione sono i vertici di un poligono di 49 lati inscritti in una circonferenza
\end{enumerate}
\end{quiz}
% SOLUZIONE: e
% SPIEGAZIONE: $49x^{48}+4>0$, quindi esiste un solo zero reale, positivo. Le radici non sono in generale equidistanti su una circonferenza.

% SIMULAZIONE / QUESITO 4 (pagina 23)
\begin{quiz}
Una primitiva della funzione $f(x)=\dfrac1{(x-5)^2}$
\begin{enumerate}[label=\alph*),leftmargin=1.3cm,itemsep=3pt,topsep=5pt]
\item $F(x)=\frac{x-4}{5-x}$, $x\in I=(5,+\infty)$
\item $F(x)=\frac{x-4}{5-x}$, $x\in I=\mathbb R\setminus\{5\}$
\item $F(x)=\frac1{x-5}+x$, $x\in I=(-\infty,5)$
\item $F(x)=\frac{x}{5-x}$, $x\in I=(5,+\infty)$
\item nessuna delle altre risposte \`e corretta
\end{enumerate}
\end{quiz}
% SOLUZIONE: a
% SPIEGAZIONE: $\frac{x-4}{5-x}=-1-\frac1{x-5}$ e la derivata \`e $1/(x-5)^2$.

% SIMULAZIONE / QUESITO 5 (pagina 23)
\begin{quiz}
La derivata della funzione $f(x)=2^{\cos(\pi x)}$ \`e:
\begin{enumerate}[label=\alph*),leftmargin=1.3cm,itemsep=3pt,topsep=5pt]
\item $-\pi\sin(\pi x)\,2^{\cos(\pi x)}$
\item $-\log2\sin(\pi x)\,2^{\cos(\pi x)}$
\item $-\pi\log2\sin(\pi x)\,2^{\cos(\pi x)}$
\item $-\pi\log2\,2^{\cos(\pi x)}$
\item $2^{\cos(\pi x)}$
\end{enumerate}
\end{quiz}
% SOLUZIONE: c
% SPIEGAZIONE: Per $a^{u(x)}$, la derivata \`e $a^{u(x)}\log(a)u\prime(x)$.

% SIMULAZIONE / QUESITO 6 (pagina 23)
\begin{quiz}
Quale delle seguenti propriet\`a NON \`e soddisfatta da $z=2+3i$?
\begin{enumerate}[label=\alph*),leftmargin=1.3cm,itemsep=3pt,topsep=5pt]
\item $z\overline z\in\mathbb R$
\item $z+\overline z\in\mathbb R$
\item $z-\overline z$ \`e un numero immaginario
\item $|z|z\in\mathbb R$
\item $z^2-6iz\in\mathbb R$
\end{enumerate}
\end{quiz}
% SOLUZIONE: d
% SPIEGAZIONE: $|z|z=\sqrt{13}(2+3i)$ non \`e reale.

\par\noindent\hfill{\color{cyan!65!black}\scriptsize Pag. originale 24/521}\par\smallskip

% SIMULAZIONE / QUESITO 7 (pagina 24)
\begin{quiz}
$\left(\frac12+\frac{\sqrt3}{2}i\right)^9=$
\begin{enumerate}[label=\alph*),leftmargin=1.3cm,itemsep=3pt,topsep=5pt]
\item $i$
\item $-i$
\item $\frac12+\frac{\sqrt3}2i$
\item $1$
\item $-1$
\end{enumerate}
\end{quiz}
% SOLUZIONE: e
% SPIEGAZIONE: La base \`e $e^{i\pi/3}$; elevando alla nona si ottiene $e^{i3\pi}=-1$.

% SIMULAZIONE / QUESITO 8 (pagina 24)
\begin{quiz}
Lo sviluppo di Mac Laurin di ordine 2 della funzione $f(x)=3^{2\sin x}$ \`e:
\begin{enumerate}[label=\alph*),leftmargin=1.3cm,itemsep=3pt,topsep=5pt]
\item $1+2x\log3+2x^2\log^23$
\item $1+2x\log3+2x^2\log3+o(x^2)$
\item $3+2x\log3+2x^2\log^23$
\item $1+2x\log3+2x^2\log^23+o(x^2)$
\item $2x\log3+2x^2\log^23+o(x^2)$
\end{enumerate}
\end{quiz}
% SOLUZIONE: d
% SPIEGAZIONE: $3^{2\sin x}=e^{2\log3(x+o(x^2))}=1+2x\log3+2x^2\log^23+o(x^2)$.

% SIMULAZIONE / QUESITO 9 (pagina 24)
\begin{quiz}
Quale delle seguenti propriet\`a \`e soddisfatta dal numero complesso $z=\dfrac{1-i}{1+i}$?
\begin{enumerate}[label=\alph*),leftmargin=1.3cm,itemsep=3pt,topsep=5pt]
\item $z=\sqrt2e^{-i\pi/4}$
\item $z^2=-2i$
\item $z^4=-4$
\item $\operatorname{Re}(z)=1$
\item $z=-i$
\end{enumerate}
\end{quiz}
% SOLUZIONE: e
% SPIEGAZIONE: Moltiplicando per $1-i$ si ottiene $z=(1-i)^2/2=-i$.

% SIMULAZIONE / QUESITO 10 (pagina 24)
\begin{quiz}
Sia data la funzione $f(x)=(x+1)(x-1)(x-4)$. Quanti punti critici ha $g(x)=\cos f(x)$?
\begin{enumerate}[label=\alph*),leftmargin=1.3cm,itemsep=3pt,topsep=5pt]
\item almeno 5
\item $5$
\item meno di 5
\item $3$
\item non \`e possibile saperlo
\end{enumerate}
\end{quiz}
% SOLUZIONE: a
% SPIEGAZIONE: $g\prime=-\sin(f)f\prime$: i tre zeri di $f$ e i due zeri distinti di $f\prime$ forniscono almeno cinque punti critici.

% SIMULAZIONE / QUESITO 11 (pagina 24)
\begin{quiz}
Se $g:\mathbb R\to\mathbb R$ \`e derivabile con $\lim_{x\to-\infty}g(x)=0$ e $\lim_{x\to+\infty}g(x)=0$, allora
\begin{enumerate}[label=\alph*),leftmargin=1.3cm,itemsep=3pt,topsep=5pt]
\item $g(x)<1$, $x<0$
\item $g(x)<1$, $x>0$
\item $\exists x_0\in\mathbb R$ tale che $g\prime(x_0)=0$
\item $f(x)>0$ per ogni $x\in\mathbb R$
\item $\exists x_0\in\mathbb R$ tale che $g(x_0)=0$
\end{enumerate}
\end{quiz}
% SOLUZIONE: c
% SPIEGAZIONE: Applicando Rolle a due punti con lo stesso valore, oppure usando un estremo interno quando $g$ non \`e costante, si trova $g\prime(x_0)=0$.

% SIMULAZIONE / QUESITO 12 (pagina 24)
\begin{quiz}
$\displaystyle\lim_{x\to0^+}\frac{\sqrt{x}-2e^x\sin x}{\log(1+x^2)-\sqrt{4x}}=$
\begin{enumerate}[label=\alph*),leftmargin=1.3cm,itemsep=3pt,topsep=5pt]
\item $-1$
\item $1$
\item $2$
\item $1/2$
\item $-1/2$
\end{enumerate}
\end{quiz}
% SOLUZIONE: e
% SPIEGAZIONE: Per $x\to0^+$, $2e^x\sin x=O(x)=o(\sqrt{x})$ e $\log(1+x^2)=O(x^2)=o(\sqrt{x})$. Inoltre $\sqrt{4x}=2\sqrt{x}$, quindi il rapporto tende a $\sqrt{x}/(-2\sqrt{x})=-1/2$.

\par\noindent\hfill{\color{cyan!65!black}\scriptsize Pag. originale 25/521}\par\smallskip

% SIMULAZIONE / QUESITO 13 (pagina 25)
\begin{quiz}
Sia data la funzione $f(x)=2x^3+3e^{2x}$. Allora $(f^{-1})\prime(3)=$
\begin{enumerate}[label=\alph*),leftmargin=1.3cm,itemsep=3pt,topsep=5pt]
\item $1/8$
\item $-1/8$
\item $1/6$
\item $-1/6$
\item $1/(54+e^6)$
\end{enumerate}
\end{quiz}
% SOLUZIONE: c
% SPIEGAZIONE: $f(0)=3$, $f\prime(0)=6$, quindi $(f^{-1})\prime(3)=1/6$.

% SIMULAZIONE / QUESITO 14 (pagina 25)
\begin{quiz}
Se $f(x)=|\cos(\pi x)|^3+2|x|$, allora $f\prime(1)=$
\begin{enumerate}[label=\alph*),leftmargin=1.3cm,itemsep=3pt,topsep=5pt]
\item $-1$
\item $1$
\item $2$
\item $3$
\item non esiste
\end{enumerate}
\end{quiz}
% SOLUZIONE: c
% SPIEGAZIONE: Vicino a $x=1$ il termine $|\cos(\pi x)|^3$ ha derivata nulla nel punto, mentre $2|x|$ ha derivata $2$.

% SIMULAZIONE / QUESITO 15 (pagina 25)
\begin{quiz}
Se $\lim_{n\to+\infty}(|a_n|+|b_n|)=0$, quale delle seguenti affermazioni \`e sicuramente vera?
\begin{enumerate}[label=\alph*),leftmargin=1.3cm,itemsep=3pt,topsep=5pt]
\item $\lim a_n=\lim b_n=0$
\item $\lim a_n=1$, $\lim b_n=-1$
\item $\lim a_n=0$, $\lim b_n=+\infty$
\item nessuna delle altre risposte \`e corretta
\item $\lim a_n=+\infty$, $\lim b_n=-\infty$
\end{enumerate}
\end{quiz}
% SOLUZIONE: a
% SPIEGAZIONE: Entrambi gli addendi sono non negativi, dunque ciascuno tende a zero.

% SIMULAZIONE / QUESITO 16 (pagina 25)
\begin{quiz}
$\displaystyle\int x(2-3x^2)^5\,dx=$
\begin{enumerate}[label=\alph*),leftmargin=1.3cm,itemsep=3pt,topsep=5pt]
\item $\frac1{36}(2-3x^2)^6+C$
\item $\frac1{36}(3x^2-2)^6+C$
\item $\frac16(2-3x^2)^6+C$
\item $-\frac1{36}(2-3x^2)^6+C$
\item $-\frac1{36}(2-3x^2)^5+C$
\end{enumerate}
\end{quiz}
% SOLUZIONE: d
% SPIEGAZIONE: Con $u=2-3x^2$, $du=-6x\,dx$, perci\`o l\'integrale \`e $-u^6/36+C$.

% SIMULAZIONE / QUESITO 17 (pagina 25)
\begin{quiz}
$f(x)=\cos x$ \`e soluzione dell\'equazione differenziale:
\begin{enumerate}[label=\alph*),leftmargin=1.3cm,itemsep=3pt,topsep=5pt]
\item $(y\prime)^2-y^2=1$
\item $y\prime\prime\prime+y=0$
\item $y\prime\prime-y=0$
\item $(y\prime)^2+y^2=1$
\item $y\prime\prime-2y=0$
\end{enumerate}
\end{quiz}
% SOLUZIONE: d
% SPIEGAZIONE: $y\prime=-\sin x$, quindi $\sin^2x+\cos^2x=1$.

\par\noindent\hfill{\color{cyan!65!black}\scriptsize Pag. originale 26/521}\par\smallskip

% SIMULAZIONE / QUESITO 18 (pagina 26)
\begin{quiz}
Quale delle seguenti funzioni NON \`e una primitiva della funzione $\dfrac{2x}{3x^2+3}$?
\begin{enumerate}[label=\alph*),leftmargin=1.3cm,itemsep=3pt,topsep=5pt]
\item $F(x)=\frac13\log(3x^2+3)$, $x\in\mathbb R$
\item $F(x)=\frac13\log(9(x^2+1))$, $x\in\mathbb R$
\item $F(x)=\frac13\log(x^2+1)+1$, $x\in\mathbb R^+$
\item $F(x)=\frac13\log(x^2+1)$, $x\in\mathbb R$
\item $F(x)=\frac13\log\left(\frac{x^2+1}{2x}\right)$, $x\in\mathbb R$
\end{enumerate}
\end{quiz}
% SOLUZIONE: e
% SPIEGAZIONE: La (e) non \`e definita su tutta $\mathbb R$ e la sua derivata include $-1/(3x)$.

% SIMULAZIONE / QUESITO 19 (pagina 26)
\begin{quiz}
Sia $f(x)$ una funzione non identicamente nulla e $f\prime(x)$ la sua derivata. Quale delle seguenti affermazioni \`e SBAGLIATA?
\begin{enumerate}[label=\alph*),leftmargin=1.3cm,itemsep=3pt,topsep=5pt]
\item $\int\frac{f\prime(x)}{f(x)}dx=\log|f(x)|$
\item $\int\frac{f\prime(x)}{f(x)}dx=2\log|f(x)|$
\item $\int\frac{f\prime(x)}{f(x)}dx=\log(e^3|f(x)|)$
\item $\int\frac{f\prime(x)}{f(x)}dx=\log(|f(x)|\sin(\pi/6))$
\item $\int\frac{f\prime(x)}{f(x)}dx=\log(2|f(x)|)$
\end{enumerate}
\end{quiz}
% SOLUZIONE: b
% SPIEGAZIONE: Una primitiva \`e $\log|f(x)|+C$ dove $f\ne0$; moltiplicare il logaritmo per $2$ ne raddoppia la derivata.

% SIMULAZIONE / QUESITO 20 (pagina 26)
\begin{quiz}
Quale delle seguenti propriet\`a \`e FALSA?
\begin{enumerate}[label=\alph*),leftmargin=1.3cm,itemsep=3pt,topsep=5pt]
\item $\sqrt[3]{-i}=e^{i(-\pi/6+2k\pi/3)}$, $k=0,1,2$
\item le tre radici cubiche di $-i$ sono vertici di un triangolo equilatero sulla circonferenza unitaria
\item una delle radici cubiche di $-i$ \`e $w=i$
\item due delle radici cubiche di $-i$ sono tra loro complesse coniugate
\item $(-\frac{\sqrt3}{2}-\frac12i)^3=-i$
\end{enumerate}
\end{quiz}
% SOLUZIONE: d
% SPIEGAZIONE: Le radici hanno argomenti $-\pi/6$, $\pi/2$, $7\pi/6$: nessuna coppia \`e coniugata.

% SIMULAZIONE / QUESITO 21 (pagina 26)
\begin{quiz}
Quale delle seguenti funzioni NON \`e una primitiva di $f(x)=1/x$?
\begin{enumerate}[label=\alph*),leftmargin=1.3cm,itemsep=3pt,topsep=5pt]
\item $F(x)=\log x$ in $I=(0,+\infty)$
\item $F(x)=\log(x/3)$ in $I=(0,+\infty)$
\item $F(x)=\log(-x)$ in $I=(-\infty,0)$
\item $F(x)=\log|x|$ in $I=\mathbb R\setminus\{0\}$
\item $F(x)=\log(-3x)$ in $I=(-\infty,0)$
\end{enumerate}
\end{quiz}
% SOLUZIONE: d
% SPIEGAZIONE: Il testo usa la definizione di primitiva su un intervallo: $\mathbb R\setminus\{0\}$ non \`e un intervallo, pur essendo la derivata $1/x$ su ciascuna componente.

% SIMULAZIONE / QUESITO 22 (pagina 26)
\begin{quiz}
Quale delle seguenti affermazioni NON \`e soddisfatta dalle due funzioni $F(x)=\arctan x$ e $G(x)=-\arctan(1/x)$?
\begin{enumerate}[label=\alph*),leftmargin=1.3cm,itemsep=3pt,topsep=5pt]
\item $F$ e $G$ sono primitive della stessa funzione sullo stesso intervallo $I$
\item $F(x)=G(x)+\pi/2$, $x\in(0,+\infty)$
\item $F(x)=G(x)-\pi/2$, $x\in(-\infty,0)$
\item $F$ e $G$ differiscono per una costante additiva che dipende dall'intervallo $I$
\item $F(x)=G(x)+\pi/2$, $x\in\mathbb R\setminus\{0\}$
\end{enumerate}
\end{quiz}
% SOLUZIONE: e
% SPIEGAZIONE: La costante \`e $\pi/2$ per $x>0$ e $-\pi/2$ per $x<0$, quindi non \`e unica su $\mathbb R\setminus\{0\}$.

\par\noindent\hfill{\color{cyan!65!black}\scriptsize Pag. originale 29/521}\par\smallskip

% SIMULAZIONE / QUESITO 1 (pagina 29)
\begin{quiz}
Siano date due funzioni $f,g:\mathbb R\to\mathbb R$ derivabili con $f(0)=g(0)$ e $f(6)=g(6)$. Quale affermazione \`e necessariamente corretta?
\begin{enumerate}[label=\alph*),leftmargin=1.3cm,itemsep=3pt,topsep=5pt]
\item $\nexists c\in(0,6):f\prime(c)=g\prime(c)=0$
\item $\exists c_1,c_2\in(0,6):f\prime(c_1)=g\prime(c_2)$
\item $\exists c\in(0,6):f\prime(c)=g\prime(c)=0$
\item le funzioni su $[0,6]$ non ammettono massimo e minimo assoluto
\item le immagini delle due funzioni su $[0,6]$ sono uguali
\end{enumerate}
\end{quiz}
% SOLUZIONE: b
% SPIEGAZIONE: Per Rolle applicato a $f-g$ esiste $c$ con $f\prime(c)=g\prime(c)$; basta prendere $c_1=c_2=c$.

% SIMULAZIONE / QUESITO 2 (pagina 29)
\begin{quiz}
Lo sviluppo di MacLaurin di ordine 4 della funzione $f(x)=\dfrac{\sqrt[3]{1+x^4}}{\sqrt[4]{1+x^3}}$ \`e:
\begin{enumerate}[label=\alph*),leftmargin=1.3cm,itemsep=3pt,topsep=5pt]
\item $1-\frac{x^2}2+\frac{x^3}3+o(x^4)$
\item $1-\frac{x^3}4+\frac{x^4}3+o(x^4)$
\item $-\frac{x^3}2+\frac{x^4}3+o(x^4)$
\item $1+\frac{x^3}2+\frac{x^4}3+o(x^4)$
\item $1-\frac{x^3}2-\frac{x^4}3+o(x^4)$
\end{enumerate}
\end{quiz}
% SOLUZIONE: b
% SPIEGAZIONE: $(1+x^4)^{1/3}=1+x^4/3+o(x^4)$, $(1+x^3)^{-1/4}=1-x^3/4+o(x^4)$.

% SIMULAZIONE / QUESITO 3 (pagina 29)
\begin{quiz}
Sia data la successione $a_n=n\sin(\frac\pi2n)$. Quale affermazione \`e vera?
\begin{enumerate}[label=\alph*),leftmargin=1.3cm,itemsep=3pt,topsep=5pt]
\item diverge a $-\infty$
\item diverge oscillando
\item converge a zero
\item $\lim_{n\to+\infty}a_n=+\infty$
\item \`e indeterminata perch'e il limite non esiste
\end{enumerate}
\end{quiz}
% SOLUZIONE: b
% SPIEGAZIONE: $\sin(\pi n/2)$ assume ciclicamente $1,0,-1,0$, quindi la successione diverge oscillando e non ha limite.

% SIMULAZIONE / QUESITO 4 (pagina 29)
\begin{quiz}
Sia data la funzione $f(x)=2^{\cos x}$. Quale affermazione \`e vera?
\begin{enumerate}[label=\alph*),leftmargin=1.3cm,itemsep=3pt,topsep=5pt]
\item la funzione non \`e periodica
\item $\lim_{x\to+\infty}f(x)=+\infty$
\item la funzione si annulla infinite volte
\item l'insieme immagine \`e $[0,+\infty)$
\item la funzione \`e limitata
\end{enumerate}
\end{quiz}
% SOLUZIONE: e
% SPIEGAZIONE: Poich\'e $-1\le\cos x\le1$, si ha $1/2\le2^{\cos x}\le2$.

% SIMULAZIONE / QUESITO 5 (pagina 29)
\begin{quiz}
Sia $f:\mathbb R\to\mathbb R$ derivabile in $0$ e $f(0)=0$. Il rapporto incrementale della funzione $f^3(x)$ in $0$ vale:
\begin{enumerate}[label=\alph*),leftmargin=1.3cm,itemsep=3pt,topsep=5pt]
\item $\bigl(\frac{f(x)-f(0)}x\bigr)^2$
\item $\frac{f(x^3)}{x^2}$
\item $\frac{f^3(x)}x$
\item $\frac{f(x^3)}x$
\item $3f^2(x)f\prime(x)$
\end{enumerate}
\end{quiz}
% SOLUZIONE: c
% SPIEGAZIONE: Per definizione $[f(x)^3-f(0)^3]/(x-0)=f(x)^3/x$.

\par\noindent\hfill{\color{cyan!65!black}\scriptsize Pag. originale 30/521}\par\smallskip

% SIMULAZIONE / QUESITO 6 (pagina 30)
\begin{quiz}
Il numero complesso $z=i^7(2-3i)^2$ vale:
\begin{enumerate}[label=\alph*),leftmargin=1.3cm,itemsep=3pt,topsep=5pt]
\item $12-7i$
\item $12+5i$
\item $-6+5i$
\item $5+6i$
\item $-12+5i$
\end{enumerate}
\end{quiz}
% SOLUZIONE: e
% SPIEGAZIONE: $(2-3i)^2=-5-12i$ e $i^7=-i$, quindi $(-i)(-5-12i)=-12+5i$. L\'opzione corretta \`e (e), non (a).

% SIMULAZIONE / QUESITO 7 (pagina 30)
\begin{quiz}
$\displaystyle\int_1^{+\infty}\frac3{x^3}\,dx=$
\begin{enumerate}[label=\alph*),leftmargin=1.3cm,itemsep=3pt,topsep=5pt]
\item $+\infty$
\item $3/2$
\item $-\infty$
\item $2/3$
\item $1/2$
\end{enumerate}
\end{quiz}
% SOLUZIONE: b
% SPIEGAZIONE: $\int_1^{+\infty}3x^{-3}dx=[-3/(2x^2)]_1^{+\infty}=3/2$.

% SIMULAZIONE / QUESITO 8 (pagina 30)
\begin{quiz}
Una delle risposte \`e FALSA, quale? $\displaystyle\int_0^{\pi/4}\tan x\,dx=$
\begin{enumerate}[label=\alph*),leftmargin=1.3cm,itemsep=3pt,topsep=5pt]
\item $-\log(\sqrt2/2)$
\item $\log(2/\sqrt2)$
\item $-\log(1/\sqrt2)$
\item $\log2$
\item $\frac12\log2$
\end{enumerate}
\end{quiz}
% SOLUZIONE: d
% SPIEGAZIONE: L\'integrale \`e $-\log(\cos\pi/4)=\frac12\log2$; (d) vale il doppio.

% SIMULAZIONE / QUESITO 9 (pagina 30)
\begin{quiz}
$\displaystyle\lim_{x\to0}\frac{\sqrt[3]{1+3x}-1}{7x+20x^4}$
\begin{enumerate}[label=\alph*),leftmargin=1.3cm,itemsep=3pt,topsep=5pt]
\item $1/7$
\item $3/7$
\item $1/21$
\item $0$
\item $+\infty$
\end{enumerate}
\end{quiz}
% SOLUZIONE: a
% SPIEGAZIONE: $\sqrt[3]{1+3x}=1+x+o(x)$, quindi il rapporto tende a $1/7$.

% SIMULAZIONE / QUESITO 10 (pagina 30)
\begin{quiz}
Se $f(x)=o(x^2)$ per $x\to0$, allora necessariamente:
\begin{enumerate}[label=\alph*),leftmargin=1.3cm,itemsep=3pt,topsep=5pt]
\item $\lim_{x\to0}f(x)/x^3=0$
\item $\lim_{x\to0}2f(x)/(x^2+3x^3)=1$
\item $\lim_{x\to0}f(x)/x=0$
\item $\lim_{x\to+\infty}f(x)/x=0$
\item $\lim_{x\to0}f(x)/\sqrt{x}=+\infty$
\end{enumerate}
\end{quiz}
% SOLUZIONE: c
% SPIEGAZIONE: $f(x)/x=[f(x)/x^2]x\to0$.
Risposta completata
